\input format.tex

%\usepackage{arydshln}

\begin{document}

%%%
\setlength{\arrayrulewidth}{0.5pt}
\color{gray2}
\vspace*{0mm}
\fontsize{9.3pt}{9.8pt}\selectfont
\begin{tctabularx}{tabularx={p{15.4cm}<{\centering}p{0pt}}}
\\[-6.769pt]
\cellcolor{kelly} \raisebox{1pt}{\color{black70}\xiaosihao\lantxh\bf 肠道调节方案} \\
\end{tctabularx}

{
\lantxh
\vspace*{-4.25mm}
\noindent
\begin{longtable}{|@{~~~}p{15.35cm}|}
\hline
\parbox[c]{\hsize}{\vskip7pt\xiaowuhao {\lantxh\bf 补充益生菌、益生元和膳食纤维}\\1.可选择性服用益生元：菊粉、低聚果糖、低聚半乳糖、低聚寡糖等。

2.每天服用适量益生菌产品：酸奶（每天1\textasciitilde 2瓶）、含乳双歧杆菌固体饮料等，帮助降低低密度脂蛋白。

3.每天选择性服用膳食纤维，这将有益于您的肠道健康。 
\vskip7pt}\\
\hline
\parbox[c]{\hsize}{\vskip7pt\xiaowuhao {\lantxh\bf 补充营养品}\\可选择性补充含有银杏提取物、绞股蓝皂甙、白藜芦醇功效成分的保健品，辅助您防治心脑血管疾病。
\vskip7pt}\\
\hline
\end{longtable}
}


\fontsize{9.3pt}{9.8pt}\selectfont
\begin{tctabularx}{tabularx={p{15.4cm}<{\centering}p{0pt}}}
\\[-6.769pt]
\cellcolor{kelly} \raisebox{1pt}{\color{black70}\xiaosihao\lantxh\bf 运动方案} \\
\end{tctabularx}

{
\lantxh
\vspace*{-4.25mm}
\noindent
\begin{longtable}{|p{1.25cm}@{}p{2cm}|p{11.732cm}|}
\hline
\parbox[c]{\hsize}{\vskip7pt\raisebox{-1.5ex}{\includegraphics[width=.9cm]{yundong.pdf}}\vskip7pt} & \parbox[c]{\hsize}{\centering\xiaowuhao\vskip7pt{运动}\vskip7pt} &
\parbox[c]{\hsize}{\vskip7pt\xiaowuhao {1.不宜进行剧烈运动。

2.应根据自身身体状况控制运动强度，避免因心率过快、血压升高而诱发心律失常、头痛、头晕、心绞痛等不适。建议选择中低强度的有氧运动，如快步走、游泳、舞蹈、自行车。配合增加肌肉量的运动，如俯卧撑、哑铃、深蹲等。

3.每次运动时间约30分钟，每周3-5次。}\vskip7pt}\\
\hline
\end{longtable}
}
\end{document}
